\chapter{Bridges To Nonintegrable Two-Dimensional QFT And CFT}
\label{ch:integrable-bridges-nonintegrable-two-dimensional-qft}

\ConventionsMixedCFT

Exact integrability is a special structure, but the machinery developed in
this volume also gives controlled coordinates near nonintegrable theories.
This chapter formulates such bridges through deformation data, broken
conservation laws, finite-volume regulators, and spectral comparison with
two-dimensional CFT.

\section{Integrability-Breaking Datum}

Let \(\mathcal Q_0\) be an integrable massive two-dimensional QFT with local
operator algebra \(\mathcal A_0\), Hilbert space \(\Hilb_0\), Hamiltonian
\[
  H_0,
\]
and commuting conserved charges
\[
  Q_s,\qquad s\in\mathcal S\subset\mathbb Z .
\]
An integrability-breaking deformation is specified by a local operator
\(\mathcal O_b\), a real coupling \(\epsilon\), and a regulator prescription
for the perturbed Hamiltonian
\[
  H(\epsilon,L,\Lambda)
  =
  H_0(L,\Lambda)
  +
  \epsilon\int_0^L dx\,\mathcal O_b(x)_{\Lambda}
  +
  \sum_a c_a(\epsilon,\Lambda)\int_0^L dx\,\mathcal O_a(x)_{\Lambda}.
\]
Here \(L\) is the spatial circumference, \(\Lambda\) is a UV cutoff, and the
last sum is the finite list of local coordinates required by the chosen
regularization.  The deformation is a QFT datum only together with these
local coordinates.

\section{Broken Conservation Laws}

Suppose \(Q_s\) is conserved in the integrable theory:
\[
  [H_0,Q_s]=0.
\]
At first order in \(\epsilon\),
\[
  \frac{\dd Q_s}{\dd t}
  =
  \ii[H(\epsilon),Q_s]
  =
  \ii\epsilon\int dx\,[\mathcal O_b(x),Q_s]+O(\epsilon^2).
\]
Equivalently, if \(Q_s=\int dx\,j^0_s(x)\), then the continuity equation
becomes
\[
  \partial_\mu j^\mu_s(x)
  =
  \epsilon\,\mathcal B_s(x)+O(\epsilon^2),
\]
where \(\mathcal B_s\) is the local operator determined by the commutator
with the perturbing density.  Matrix elements of \(\mathcal B_s\) between
asymptotic integrable states diagnose which multiparticle processes are
created by the deformation.

\section{Form-Factor Perturbation}

Let
\[
  \ket{\theta_1,a_1;\ldots;\theta_n,a_n}_{\mathrm{in}}
\]
be an \(n\)-particle incoming state of the integrable theory.  The first
order transition matrix element generated by \(\mathcal O_b\) is
\[
  \mathcal M_{\beta\alpha}^{(1)}
  =
  \epsilon\int d^2x\,
  {}_{\mathrm{out}}\!\bra{\beta}\mathcal O_b(x)\ket{\alpha}_{\mathrm{in}}.
\]
Translation covariance gives
\[
  {}_{\mathrm{out}}\!\bra{\beta}\mathcal O_b(x)\ket{\alpha}_{\mathrm{in}}
  =
  \ee^{\ii(P_\beta-P_\alpha)\cdot x}
  {}_{\mathrm{out}}\!\bra{\beta}\mathcal O_b(0)\ket{\alpha}_{\mathrm{in}}.
\]
Therefore the infinite-volume expression contains the energy-momentum
distribution
\[
  (2\pi)^2\delta^{(2)}(P_\beta-P_\alpha)
  \epsilon\,{}_{\mathrm{out}}\!\bra{\beta}\mathcal O_b(0)\ket{\alpha}_{\mathrm{in}}.
\]
The local matrix element is an analytic continuation of the integrable
form-factor data, with disconnected pieces and kinematic poles treated by the
crossing prescription of the form-factor chapter.

\section{Widths From Finite-Volume Perturbation}

Let \(\ket{A,L}\) be a one-particle finite-volume eigenstate of \(H_0(L)\),
and let \(\ket{n,L}\) be finite-volume multiparticle states with energies
\(E_n(L)\).  Second-order time-dependent perturbation theory gives the
transition probability
\[
  P_{A\to n}(T,L)
  =
  \epsilon^2
  \left|
    \int_0^T dt\,\ee^{\ii(E_n-E_A)t}
    \bra{n,L}V_b\ket{A,L}
  \right|^2
  +
  O(\epsilon^3),
\]
where \(V_b=\int_0^L dx\,\mathcal O_b(x)\).  Since
\[
  \left|\int_0^Tdt\,\ee^{\ii\Delta E t}\right|^2
  =
  \frac{4\sin^2(\Delta E T/2)}{(\Delta E)^2},
\]
the large-\(T\) distributional limit is
\[
  \frac1T
  \left|\int_0^Tdt\,\ee^{\ii\Delta E t}\right|^2
  \longrightarrow
  2\pi\delta(\Delta E).
\]
After the \(L\to\infty\) density of states is inserted, the decay width is
\[
  \Gamma_A
  =
  2\pi\epsilon^2
  \sum_\beta
  \rho_\beta(E_A)
  \left|
    {}_{\mathrm{out}}\!\bra{\beta}\mathcal O_b(0)\ket{A}_{\mathrm{in}}
  \right|^2,
\]
within the stated perturbative and infrared hypotheses.

\section{Truncated Conformal Space As Regulator}

Let \(\mathcal C\) be a two-dimensional CFT on a circle of circumference
\(L\).  The CFT Hamiltonian is
\[
  H_{\mathrm{CFT}}
  =
  \frac{2\pi}{L}
  \left(L_0+\bar L_0-\frac{c}{12}\right).
\]
Perturbing by relevant operators gives
\[
  H=H_{\mathrm{CFT}}
  +\sum_i \lambda_i
  \left(\frac{L}{2\pi}\right)^{1-\Delta_i}
  \int_0^{2\pi}d\theta\,\mathcal O_i(\theta).
\]
A truncation cutoff \(E_{\mathrm{cut}}\) keeps only CFT states below that
energy.  Matrix elements are computed from CFT three-point coefficients and
descendant Ward identities.  Counterterms depend on the discarded high-energy
states and must be included as local coordinates in the truncated Hamiltonian.

\section{Comparison Objects}

A bridge between integrable and nonintegrable two-dimensional QFT compares
objects, not names:
\[
\begin{gathered}
  \text{finite-volume spectrum},\qquad
  \text{local-operator matrix elements},\\
  \text{resonance poles},\qquad
  \text{hydrodynamic or thermal late-time behavior}.
\end{gathered}
\]
An integrable point may control nearby data through perturbation theory in
\(\epsilon\), while a CFT point may control nearby data through a relevant
deformation.  The two expansions meet only after the finite-volume regulator,
operator normalization, and continuation to infinite volume are matched.

\begin{openproblem}[Nonintegrable deformation]
\label{op:controlled-nonintegrable-deformation}
Construct a nonintegrable two-dimensional QFT near an integrable point with a
finite-volume regulator, local counterterm coordinates, form-factor
perturbation theory, resonance pole analysis, and a proof of convergence to a
local continuum QFT in a specified scaling domain.
\end{openproblem}
